\begin{lemma}
\label{editorial-source-lemma-polynomials-divide}
Let $K$ be a field. Let $n, m \in \mathbf{N}$ and
$a_0, \ldots, a_{n - 1}, b_0, \ldots, b_{m - 1} \in K$.
If the polynomial $x^n + a_{n - 1}x^{n - 1} + \ldots + a_0$
divides the polynomial $x^m + b_{m - 1} x^{m - 1} + \ldots + b_0$
in $K[x]$ then
\begin{enumerate}
\item $a_0, \ldots, a_{n - 1}$ are integral over any subring
$R_0$ of $K$ containing the elements $b_0, \ldots, b_{m - 1}$, and
\item each $a_i$ lies in $\sqrt{(b_0, \ldots, b_{m-1})R}$
for any subring $R \subset K$ containing the elements
$a_0, \ldots, a_{n - 1}, b_0, \ldots, b_{m - 1}$.
\end{enumerate}
\end{lemma}

\begin{proof}
Write
$$
Q=x^n+\sum_{j=0}^{n-1}a_jx^j,\qquad
P=x^m+\sum_{i=0}^{m-1}b_ix^i.
$$
Since both are monic and $Q$ divides $P$ over the field $K$,
we have $m\geq n$. Choose a splitting field $L/K$ and retain
the roots with their multiplicities:
$$
P=\prod_{i=1}^m(x-\beta_i),\qquad \beta_i\in L.
$$
See Fields, Section \ref{fields-section-splitting-fieds}.
Unique factorization in $L[x]$ gives an indexed subset
$J\subseteq\{1,\ldots,m\}$ of cardinality $n$ with
$Q=\prod_{i\in J}(x-\beta_i)$. Equal roots remain distinct
indices in this expression. Every $\beta_i$ is integral over
$R_0$ by the original $P$, and every original coefficient is
$$
a_j=(-1)^{n-j}
\sum_{\substack{E\subseteq J\\ |E|=n-j}}
\prod_{i\in E}\beta_i\qquad(0\leq j<n).
$$
The subring property in Lemma \ref{editorial-source-lemma-integral-closure-is-ring}
proves (1), retaining every sign and repeated-root contribution.

For (2), first consider $m=n$. The monic quotient has degree
zero and is $1$, so $a_i=b_i$ for every $i$, which gives the
assertion directly. Now let $e=m-n>0$ and write the original
monic quotient as
$$
C=x^e+\sum_{k=0}^{e-1}c_kx^k,\qquad P=QC.
$$
Put $a_n=c_e=1$. Comparing the coefficient of $x^{n+k}$
for $k=e-1,e-2,\ldots,0$ gives the exact recurrence
$$
c_k=b_{n+k}-
\sum_{i=\max(0,n+k-e)}^{n-1}a_i c_{n+k-i}.
$$
Every index $n+k-i$ in this sum is between $k+1$ and $e$.
Thus descending induction shows $c_k\in R$ for all $k$,
without negative coefficient indices or assumptions on $n>1$.
Set $H=\sqrt{(b_0,\ldots,b_{m-1})R}$ and
$\overline R=R/H$. This quotient is reduced. If it is the
zero ring, the required membership in $H=R$ already holds.
Otherwise the complete coefficient image of $P=QC$ is
$$
x^m+\sum_{i=0}^{m-1}\overline b_i x^i=x^m
=\left(x^n+\sum_{j=0}^{n-1}\overline a_jx^j\right)
 \left(x^e+\sum_{k=0}^{e-1}\overline c_kx^k\right).
$$
At a minimal prime $\mathfrak p$ of $\overline R$, the ring
$\overline R_{\mathfrak p}$ is a field by
Lemma \ref{algebra-lemma-minimal-prime-reduced-ring}. In its polynomial
ring unique factorization makes the first monic factor
$f=x^n+\sum_{j=0}^{n-1}\overline a_jx^j$ equal to $x^n$:
every irreducible factor divides $x$, its degree is $n$, and
its leading coefficient fixes the unit factor to be $1$.
For each $j$ choose $s_j\notin\mathfrak p$ with
$s_j\overline a_j=0$, using the equality in this localization,
and take $s=\prod_{j=0}^{n-1}s_j$.
Then there exists an $s\in\overline R$, $s\notin\mathfrak p$,
such that $s(f-x^n)=0$, with this explicit common choice.
Primality gives $\overline a_j\in\mathfrak p$ for every $j$.
As this holds at every minimal prime and their intersection
is zero in the reduced ring, every $\overline a_j$ is zero
(Lemma \ref{algebra-lemma-reduced-ring-sub-product-fields}).
Hence every $a_j$ belongs to $H$, proving (2).
\end{proof}

\begin{lemma}
\label{editorial-source-lemma-weighted-monic-factor-coefficients}
Let $R\subseteq K$ be a subring of a field and let $I\subseteq R$
be an ideal. Suppose
$$
P=x^m+\sum_{i=0}^{m-1}b_ix^i,\qquad b_i\in I^{m-i},
\qquad Q=x^n+\sum_{j=0}^{n-1}a_jx^j\in K[x]
$$
are monic, $1\leq n\leq m$, and $Q$ divides $P$.
For each $0\leq j<n$, the coefficient $a_j$ is integral over
$I^{n-j}$. More precisely, with $N=\binom{m}{n}$, it satisfies
an equation
$$
a_j^N+\sum_{\ell=0}^{N-1}C_{j,\ell}a_j^\ell=0,
\qquad C_{j,\ell}\in I^{(n-j)(N-\ell)}.
$$
\end{lemma}

\begin{proof}
In a splitting field $L/K$ retain the original indexed roots,
including multiplicities, so
$P=\prod_{r=1}^m(x-\beta_r)$. As in the proof of
Lemma \ref{editorial-source-lemma-polynomials-divide}, there is an indexed subset
$J\subseteq\{1,\ldots,m\}$ of cardinality $n$ with
$Q=\prod_{r\in J}(x-\beta_r)$. Fix $j$ and put $w=n-j$.
Let $\mathcal T$ be the set of all $n$-element subsets of
$\{1,\ldots,m\}$, so $|\mathcal T|=N$. For each $T\in\mathcal T$
define the full coefficient
$$
r_{j,T}=(-1)^w
\sum_{\substack{U\subseteq T\\|U|=w}}\prod_{r\in U}\beta_r.
$$
In particular $a_j=r_{j,J}$. Form
$$
H_j(Y)=\prod_{T\in\mathcal T}(Y-r_{j,T})
=Y^N+\sum_{\ell=0}^{N-1}C_{j,\ell}Y^\ell,
$$
where every coefficient, sign and repeated contribution is
$$
C_{j,\ell}=(-1)^{N-\ell}
\sum_{\substack{\mathcal A\subseteq\mathcal T\\|\mathcal A|=N-\ell}}
\prod_{T\in\mathcal A}r_{j,T}.
$$
Then $H_j(a_j)=0$.

Replace the $\beta_r$ temporarily by independent variables
$Z_r$. The same expression for $C_{j,\ell}$ is an integer
polynomial, symmetric in the $m$ variables and homogeneous of
total degree $D=w(N-\ell)$. It is an integer polynomial in
the elementary symmetric polynomials $e_1,\ldots,e_m$, each
monomial having weighted degree $D$ when $e_s$ has weight $s$.
Here is a proof retaining the weights. In lexicographic order
$Z_1>\cdots>Z_m$, the leading exponent vector
$(\lambda_1,\ldots,\lambda_m)$ of a nonzero symmetric
homogeneous polynomial satisfies
$\lambda_1\geq\cdots\geq\lambda_m$: an inversion would give
a larger monomial after interchanging the corresponding
variables, contradicting maximality. With $\lambda_{m+1}=0$,
the polynomial
$$
\prod_{s=1}^m e_s^{\lambda_s-\lambda_{s+1}}
$$
has that same leading monomial with coefficient $1$ and
weighted degree
$\sum_{s=1}^m s(\lambda_s-\lambda_{s+1})
=\sum_{s=1}^m\lambda_s=D$.
Subtract the integer multiple matching the leading coefficient.
The result remains symmetric and homogeneous of degree $D$,
with a smaller leading monomial. Only finitely many monomials
have total degree $D$, so this procedure terminates. It proves
an expression with integers $\kappa_{j,\ell,\nu}$:
$$
C_{j,\ell}
=\sum_{\substack{\nu\in\mathbf Z_{\geq0}^m\\
                  \sum_{s=1}^m s\nu_s=D}}
\kappa_{j,\ell,\nu}
\prod_{s=1}^m e_s(\beta_1,\ldots,\beta_m)^{\nu_s}.
$$
The original coefficients of $P$ give
$e_s(\beta_1,\ldots,\beta_m)=(-1)^s b_{m-s}$, so the
complete substituted expression is
$$
C_{j,\ell}
=\sum_{\substack{\nu\in\mathbf Z_{\geq0}^m\\
                  \sum_{s=1}^m s\nu_s=D}}
\kappa_{j,\ell,\nu}
\prod_{s=1}^m\bigl((-1)^s b_{m-s}\bigr)^{\nu_s}.
$$
Each product belongs to
$\prod_s(I^s)^{\nu_s}=I^D=I^{w(N-\ell)}$.
This proves the required equation over $R$. All identities
were proved over the integers before evaluation, so no
separability or characteristic restriction is present.

The exact Rees comparison gives a second proof of the
integrality conclusion. Put
$A=\bigoplus_{h\geq0}I^ht^h\subseteq R[t]$ and map it to
$L[t]$ by the given inclusion. Every $t\beta_r$ satisfies
$$
(t\beta_r)^m+
\sum_{i=0}^{m-1}(b_it^{m-i})(t\beta_r)^i=t^mP(\beta_r)=0.
$$
Its coefficients belong to $A$, so all these elements are
integral over $A$. The original signed identity
$$
a_jt^{n-j}=(-1)^{n-j}
\sum_{\substack{U\subseteq J\\|U|=n-j}}
\prod_{r\in U}(t\beta_r)
$$
and Lemma \ref{editorial-source-lemma-integral-closure-is-ring} show that this
element is integral over $A$. The weight-$n-j$ assertion of
Lemma \ref{editorial-source-lemma-characterize-integral-ideal} gives integrality
of the original $a_j$ over $I^{n-j}$.
\end{proof}

\begin{lemma}
\label{editorial-source-lemma-integral-ideal-minimal-polynomial-criterion}
Let $R$ be a domain, $K=\operatorname{Frac}(R)$, $I\subseteq R$
an ideal, and $g$ an algebraic element in a field extension
$L/K$. Write its monic minimal polynomial as
$Q=x^n+\sum_{j=0}^{n-1}a_jx^j\in K[x]$.
Then $g$ is integral over $I$ if and only if every $a_j$ is
integral over $I^{n-j}$. If $R$ is normal, this is equivalent
to $a_j\in\overline{I^{n-j}}\subseteq R$ for every $j$,
where the bar denotes integral closure of the ideal inside $R$.
\end{lemma}

\begin{proof}
If $g$ is integral over $I$, retain a witnessing equation
$P(g)=0$ with $P=x^m+\sum_{i=0}^{m-1}b_ix^i$ and
$b_i\in I^{m-i}$. Division by $Q$ in $K[x]$ gives a
remainder of degree less than $n$ that vanishes at $g$, hence
is zero by minimality. Thus $Q\mid P$, and
Lemma \ref{editorial-source-lemma-weighted-monic-factor-coefficients} gives
the asserted coefficient integrality.

Conversely suppose every coefficient has that integrality.
Use $A=\bigoplus_{h\geq0}I^ht^h\subseteq L[t]$ with the
original coefficient inclusion. By the weighted Rees
characterization every $h_j=a_jt^{n-j}$ is integral over $A$.
The algebra $B=A[h_0,\ldots,h_{n-1}]$ is finite over $A$:
if the chosen monic equation for $h_j$ has degree $d_j$,
the products $\prod_jh_j^{e_j}$ with $0\leq e_j<d_j$
generate it as an $A$-module, using those original monic
equations to reduce each higher power. The element $gt$
satisfies the full monic equation
$$
(gt)^n+\sum_{j=0}^{n-1}h_j(gt)^j=t^nQ(g)=0.
$$
Thus $B[gt]$ is generated over $A$ by all
$\bigl(\prod_jh_j^{e_j}\bigr)(gt)^r$ with $0\leq r<n$.
To see directly that this gives integrality over $A$, list
these generators as $v_1,\ldots,v_M$, including $1$.
Multiplication by $gt$ has relations
$(gt)v_i=\sum_k d_{ik}v_k$ with $d_{ik}\in A$.
Multiplying the vector equation by the adjugate of
$XI-(d_{ik})$ at $X=gt$ shows that the monic polynomial
$\det(XI-(d_{ik}))$ annihilates every $v_i$, including $1$.
It is therefore an integral equation for $gt$ over $A$.
The weight-one Rees characterization makes $g$ integral over
$I$, proving the converse. This argument does not require
the displayed module generators to be linearly independent.

If $R$ is normal, each integral equation for $a_j$ over an
ideal is in particular an integral equation over $R$.
Since $a_j\in K$, integral closedness of the normal domain
gives $a_j\in R$. The definition of integral closure of
$I^{n-j}$ inside $R$ now gives exactly the stated equivalence.
\end{proof}

\begin{example}
\label{editorial-source-example-weighted-coefficients-need-integral-closure}
Normality does not place these coefficients in the ideal
powers themselves. Let $k$ be any field,
$R=k[u,v]$, $K=k(u,v)$ and $I=(u^2,v^2)$.
The polynomial ring $R$ is a normal domain by
Lemma \ref{editorial-source-lemma-polynomial-domain-normal}, applied twice to $k$.
Choose a root $g$ of $x^2-u^3v$ in a field extension of $K$.
This quadratic is irreducible: in $k(u)[v]$ the exponent of
the factor $v$ defines an integer order on nonzero rational
functions by numerator order minus denominator order.
Multiplication adds this order, as follows by removing the
largest powers of $v$ and observing that the remaining
polynomials have nonzero values at $v=0$. Thus every nonzero
square in $K$ has even order, whereas $u^3v$ has order $1$.
It cannot be a square, so the quadratic has no root in $K$
and is irreducible, including in characteristic two.
Consequently $Q=x^2-u^3v$ is the minimal polynomial of $g$,
and $S=R[g]$ in the chosen extension field is a domain.
Yet the original equation
$$
g^4-u^6v^2=0,\qquad u^6v^2=(u^2)^3v^2\in I^4
$$
makes $g$ integral over $I$; its other three coefficients
are zero and belong to the required powers. The constant
coefficient $a_0=-u^3v$ of $Q$ is not in
$I^2=(u^4,u^2v^2,v^4)$: the monomial $u^3v$ is divisible
by none of these generators, and every polynomial multiple
of a generator is a sum of divisible monomials. Its distinct
monomial coefficient therefore cannot be obtained in their
sum. However
$$
a_0^2-u^6v^2=0,\qquad -u^6v^2\in I^4=(I^2)^2,
$$
so $a_0\in\overline{I^2}\setminus I^2$.
Even weight one need not give actual membership: $h=uv\in R$
satisfies $h^2-u^2v^2=0$ with $u^2v^2\in I^2$, but the
constant coefficient $-uv$ of its minimal polynomial $x-uv$
is outside $I$ by the same monomial argument.
\end{example}

\begin{lemma}
\label{editorial-source-lemma-integrally-closed-powers-monic-factors}
For a normal domain $R$, $K=\operatorname{Frac}(R)$ and an
ideal $I\subseteq R$, the following are equivalent:
\begin{enumerate}
\item Every positive power $I^w$ is integrally closed in $R$.
\item For every monic $P=x^m+\sum_{i=0}^{m-1}b_ix^i$ with
$b_i\in I^{m-i}$, every monic divisor
$Q=x^n+\sum_{j=0}^{n-1}a_jx^j\in K[x]$ of positive degree
has $a_j\in I^{n-j}$ for every $j$.
\item Every element integral over $I$ in a field extension of
$K$ has minimal-polynomial coefficients $a_j\in I^{n-j}$,
where $n$ is its minimal-polynomial degree.
\end{enumerate}
\end{lemma}

\begin{proof}
The weighted coefficient lemma gives (1)$\Rightarrow$(2):
its equations first give $a_j\in R$ by normality and then
$a_j\in I^{n-j}$ by integral closedness of that power.
Minimal-polynomial division as above gives (2)$\Rightarrow$(3).
For (3)$\Rightarrow$(2), factor the given $Q$ in $K[x]$
into monic irreducible factors, repeated with their original
multiplicities. A root of each factor in a splitting field
is a root of $P$, hence is integral over $I$. Condition (3)
gives the required powers for each irreducible factor.
Write the factors as
$Q_r=\sum_{j=0}^{d_r}c_{r,j}x^j$ for $1\leq r\leq q$,
with $c_{r,d_r}=1\in I^0$ and
$c_{r,j}\in I^{d_r-j}$ for $j<d_r$.
The full coefficient of $x^j$ in their product is
$$
\sum_{\substack{0\leq j_r\leq d_r\ (1\leq r\leq q)\\
                 j_1+\cdots+j_q=j}}
\prod_{r=1}^q c_{r,j_r}
\ \in I^{\sum_{r=1}^qd_r-j}=I^{n-j}.
$$
This retains the multiplicities and proves (2).

Finally assume (2), fix $w\geq1$, and take
$c\in\overline{I^w}\subseteq R$. Keep its defining equation
$$
c^d+\sum_{\ell=0}^{d-1}r_\ell c^\ell=0,
\qquad r_\ell\in I^{w(d-\ell)}.
$$
Set $F(Y)=Y^d+\sum_{\ell=0}^{d-1}r_\ell Y^\ell$.
The complete difference identity is
$$
F(Y)-F(c)=(Y-c)\left(
\sum_{k=0}^{d-1}Y^{d-1-k}c^k+
\sum_{\ell=1}^{d-1}r_\ell
   \sum_{k=0}^{\ell-1}Y^{\ell-1-k}c^k\right).
$$
Since $F(c)=0$, substituting $Y=x^w$ proves that $x^w-c$
divides
$P=x^{wd}+\sum_{\ell=0}^{d-1}r_\ell x^{w\ell}$.
The coefficient at $x^{w\ell}$ belongs to $I^{wd-w\ell}$;
every omitted coefficient is zero and is retained in its
required power. Condition (2) applied to this divisor gives
$-c\in I^w$, hence $c\in I^w$. The reverse inclusion
$I^w\subseteq\overline{I^w}$ follows from the degree-one
equation for each of its elements. This proves (1).
\end{proof}

\begin{remark}
\label{editorial-source-remark-weighted-coefficient-consequences}
Under the normal-domain hypotheses, the original coefficient
$a_j\in\overline{I^{n-j}}$ lies in $\sqrt I$: reduce its
defining equation modulo $\sqrt I$, where all nonleading
coefficients vanish; its positive leading power is zero, so
it belongs to that radical. If $I$ is integrally closed,
the same defining coefficients belong to $I^{N-\ell}$
because $(n-j)(N-\ell)\geq N-\ell$, and therefore $a_j$
is integral over $I$ and belongs to $I$. If $I$ is radical,
the radical conclusion already gives membership in $I$.
These claims do not assert membership in $I^{n-j}$ without
the corresponding integral-closedness condition.

The exponent $n-j$ cannot be increased in general.
For $R=k[u]$, $I=(u)$ and any $w\geq1$, take
$P=Q=x^w-u^w$. Its constant coefficient $a_0=-u^w$
belongs to $I^w$. If it were integral over $I^{w+1}$,
an equation of degree $d\geq1$ would have leading term
$(-u^w)^d$ of $u$-order $wd$, whereas each other term
$c_\ell(-u^w)^\ell$ would have order at least
$(w+1)(d-\ell)+w\ell=wd+d-\ell>wd$.
The coefficient of $u^{wd}$ in their sum would remain
$(-1)^d\ne0$, a contradiction in every characteristic.
\end{remark}